\documentclass[reqno,12pt]{amsart}
\usepackage{amsmath,amssymb,amsthm,mathtools}
\usepackage{mathrsfs}
\usepackage{bm}
\usepackage{tikz-cd}
\usepackage{booktabs}
\usepackage{xcolor}
\usepackage{hyperref}
\usepackage{float}
\usepackage{orcidlink}
\usepackage{comment}
\hypersetup{
    colorlinks=true,
    linkcolor=blue,
    citecolor=blue,
    urlcolor=blue}
\usepackage[a4paper,margin=1in]{geometry}
\newtheorem{theorem}{Theorem}[section]
\newtheorem{lemma}[theorem]{Lemma}
\newtheorem{proposition}[theorem]{Proposition}
\newtheorem{corollary}[theorem]{Corollary}
\newtheorem{assumption}{Assumption}
\theoremstyle{definition}
\newtheorem{thm}{Theorem}
\renewcommand*{\thethm}{\Alph{thm}}
\newtheorem{definition}[theorem]{Definition}
\newtheorem{remark}[theorem]{Remark}
\newtheorem{example}[theorem]{Example}
\allowdisplaybreaks
\newcommand{\F}{\mathbb{F}}
\newcommand{\Z}{\mathbb{Z}}
\newcommand{\Q}{\mathbb{Q}} 
\newcommand{\QQ}{\overline{\mathbb{Q}}}
\newcommand{\N}{\mathbb{N}}    
\newcommand{\R}{\mathbb{R}}    
\newcommand{\C}{\mathbb{C}}    
\newcommand{\FF}{{\mathcal{F}}}
\newcommand{\p}{\mathfrak{p}}
\newcommand{\q}{\mathfrak{q}}
\newcommand{\m}{\mathfrak{m}}
\newcommand{\OO}{\mathcal{O}}
\newcommand{\A}{\mathcal{A}}
\newcommand{\T}{\mathcal{T}}
\newcommand{\Gal}{\mathrm{Gal}}
\newcommand{\Hom}{\mathrm{Hom}}
\newcommand{\Ker}{\mathrm{Ker}}
\newcommand{\Coker}{\mathrm{Coker}
}
\newcommand{\ord}{\mathrm{ord}}
\newcommand{\lra}{\longrightarrow}
\newcommand{\inj}{\hookrightarrow}
\newcommand{\surj}{\twoheadrightarrow}

\pagestyle{plain}
\raggedbottom

\title{On the Diophantine Equation $\displaystyle p^x+ (2p+1)^y =z^2$ with Consecutive Exponents}

\author{Subhasis Panda} %\, \orcidlink{0009-0005-7258-4580}}

\address{Department of Mathematics, School of Advanced Sciences, VIT-AP University, Amravati 522241, Andhra Pradesh, India}

\email{subhasispanda559@gmail.com}

\subjclass[2020]{Primary 11D61; Secondary 11D41, 11A41.}

\keywords{Exponential Diophantine equation, Lifting-the-exponent lemma.}


\begin{document}


\begin{abstract} { We study the Diophantine equation $p^x+ (2p+1)^y =z^2$ over positive integers $x$, $y$ and $z$ for every odd prime $p$. We prove that $(x,y,z)=(2,1,p+1)$ is the unique solution except possibly when $2p+1$ is composite. In that case, it reduces to a family depending on one parameter, and we show that the parameter must be odd and satisfies an explicit upper bound, reducing the problem to finitely many cases.
}\end{abstract}

\maketitle

\setcounter{tocdepth}{1}
\tableofcontents
\section{Introduction}
Diophantine equations of the form $a^x+b^y=z^2$ and $a^x-b^y=z^2$,
where $a$ and $b$ are fixed positive integers, have received a lot of attention \cite{BorahDutta2022,Le1995,MinaBacani2021,ZhangLi2024,Sroysang2012,GahaMezroui2026,ChoudhurySarma2026}. In most cases, the integers a and b are fixed, and the aim is to determine all positive integer solutions $(x,y,z)$ of the corresponding equation.

For example, Sroysang\cite{Sroysang2012,Sroysang2013} studied the equations $3^x + 5^y = z^2$ and $7^x + 8^y = z^2$. 
Suvarnamani \cite{Suvarnamani2014} studied the equation $p^x+(p+1)^y=z^2$, where $p$ is an odd prime, and proved that it has a unique positive integer solution when $p=3$. Subsequently, Choudhury-Sarma \cite{ChoudhurySarma2026} investigated the  equations $(p+a)^x+p^y=z^{2}$ and $(p+a)^x-p^y=z^{2}$, where $p$ and $p+a$ are odd primes.
Later, Mina and Bacani \cite{MinaBacani2021} considered the more general family $p^x+(p+4k)^y=z^2$, where both $p$ and $p+4k$ are prime, and established several results concerning the existence of positive integer solutions. Further, Zhang and Li \cite{ZhangLi2024} investigated a general class of Diophantine equations of the form $(-1)^\alpha p^x+(-1)^\beta\left(2^k(2p-1)\right)^y=z^2$, where $(p,2p-1)$ is a prime pair. Recently, Gaha and Mezroui \cite{GahaMezroui2026} studied the equation $43^x+87^y=z^2$, under the condition that $x$ and $y$ are consecutive integers. In this paper, we investigate the general family
\begin{equation} \label{equation1}
p^x + \left ( 2p +1 \right)^y = z^2    
\end{equation}
where $p$ is an odd prime and $|x-y|=1$, without assuming that $2p+1$ is prime. 
The condition $|x-y|=1$ naturally divides the problem into two cases according to whether $x=y+1$ or $y=x+1$.  
We show that $(x,y,z)=(2,1,p+1)$ is always a solution of the equation. 
Furthermore, we prove that any remaining solution can occur only when 2p+1 is composite, in which case the original equation reduces to an auxiliary equation. Our main result is the following.

\noindent \begin{thm} (Theorem~\ref{T3.1})
\label{A}
Let $p$ be an odd prime and suppose that
\begin{equation*}
p^x + \left ( 2p +1 \right)^y = z^2,
\end{equation*}
where $x$, $y$, $z$ are positive integers with $|x-y|=1$. Then exactly one of the following holds.
\begin{enumerate}
    \item[(i)] The solution is $(x, y, z) = (2, 1, p + 1)$.
    \item[(ii)] The integer $2p + 1$ is composite and admits a coprime factorization 
    \begin{equation*}
        2p + 1 = (2k - 1)(2k + 1),
    \end{equation*}
    where $p=2k^2-1$ and $(2k + 1)^{2b+1} - (2k - 1)^{2b+1} = 2(2k^2 - 1)^a$, with $b  < \lceil \log_2 (2k^2-1) \rceil$ leaving only a finite logarithmic range of exponents for each fixed prime.
\end{enumerate}
\end{thm}
This reduction sharply limits the search for any exceptional solutions.
\section{Preliminary Results}
We prove the theorems using elementary number theory. In this section, we first state the Lifting-the-Exponent Lemma (LTE), which will be useful in our proofs. We then establish several preliminary results that will be used in the proofs of our main theorems.
\begin{lemma}  \label{lemma3}
\normalfont
Let $p$ be an odd prime, and let $a,b \in \Z$ satisfy $p | (a-b), \,\, p \nmid ab$. Then
\begin{equation*}
 v_p(a^n-b^n)=v_p(a-b)+v_p(n).   
\end{equation*}
\end{lemma}
In this paper, we discuss the above lemma with $a=q=2p+1$ and $b=1$. Since $v_p(q-1) = v_q(2q)=1$, we immediately obtain the following consequence.
\begin{corollary} \label{corollary1}
\normalfont For every integer $b \geq 0$, we have 
\begin{equation*}
v_p\!\left((2p+1)^{2b+1}-1\right) = 1+v_p(2b+1).
\end{equation*}
\end{corollary}
\begin{proof}
We apply Lemma~\ref{lemma3} with $a=2p+1, \, b=1,\, n=2b+1$. Since $v_p\!\left((2p+1)-1\right)=v_p(2p)=1$, the result follows immediately.
\end{proof}
Throughout the paper, let $p$ be an odd prime, $q:=2p+1$, and
$x,y,z$ be positive integers that satisfy $\displaystyle p^{x}+q^{y}=z^{2}$ with $|x-y|=1.$ 
Since $x$ and $y$ are consecutive integers, they have opposite parity. Therefore, exactly one of $x$ and $y$ is even and the other is odd. 
\subsection{Case A: x odd, y even}
In this case, we write $x=2a+1$ and $y=2b$, where $a,b\ge0$. The condition \(|x-y|=1\) then gives rise to exactly two sub-cases: $y=x-1$, which is equivalent to $b=a$, and $y=x+1$, which is equivalent to $b=a+1$. 
In either sub-case, the equation $p^x+q^y=z^2$ takes the form $p^{2a+1}+q^{2b}=z^2,$ or equivalently, $p^{2a+1}=z^2-q^{2b}=(z-q^b)(z+q^b)$. We first show that the two factors on the left-hand side are coprime. 
Let $d=\gcd(z-q^b,\,z+q^b)$. Then $d\mid\bigl((z+q^b)-(z-q^b)\bigr)=2q^b$ and $d\mid p^{2a+1}$. 
Since $p$ is odd, $d$ is also odd. Hence $d\mid q^b$. Therefore,
$d\mid p^{2a+1} \,\, \text{and} \,\, d\mid q^b$. Since $\gcd(p,q)=1$, we conclude that $d=1$. 
Thus, $\gcd(z-q^b,\,z+q^b)=1$. Since $z^2-q^{2b}=(z-q^b)(z+q^b)$ is a coprime factorization of a prime power, one of the two factors must be $1$.
As $z+q^b > z-q^b >0$ it follows that $z-q^b=1$ and $z+q^b=p^{2a+1}$. 
By subtracting the two equations, we obtain the following. 
\begin{equation}\label{equation2}
 p^{2a+1}-1=2q^b   
\end{equation}
First, we consider the sub-case of $y=x+1$, (equivalently, $b=a+1$), which gives us the equation $p^{2a+1}-1=2q^{a+1}$. The following Lemma shows that this equation has no solution.
\begin{lemma} \label{lemma1}
\normalfont Let $p$ be an odd prime, and $q=2p+1$. Then equation 
\begin{equation*}
p^{2a+1}-1=2q^{a+1}  
\end{equation*}
 has no solution for $a  \geq 0$.
\end{lemma}
\begin{proof}
Suppose, to the contrary, that $p^{2a+1}-1=2q^{a+1}$. Now, reducing the above equation modulo $p$, we obtain $-1\equiv 2q^{a+1} \equiv 2 \pmod p$. Thus $p\mid3$. 
As $p$ is an odd prime, we must have $p=3$, and consequently $q=7$. The equation becomes $3^{2a+1}-1=2\cdot7^{a+1}$. 
Now, reducing this equation modulo $7$, we obtain $3^{2a+1}-1 \equiv 0 \pmod 7$.
Since the multiplicative order of $3$ modulo $7$ is $6$, this gives $6 | (2a+1)$. This is impossible because $2a+1$ is odd. Hence no such $a$ exists.
\end{proof}
Next, we consider the sub-case $(y=x-1)$, (equivalently, $b=a$), with $a \ge 1$.
In this case, the equation~(\ref{equation2}) becomes $p^{2a+1}-1=2q^{a}$.
We show that the equation~(\ref{equation2}) has no solution for this sub-case as well.
\begin{lemma} \label{lemma2}
\normalfont
Let $p$ be an odd prime and let $q=2p+1$. Then the equation
\begin{equation*}
 p^{2a+1}-1=2q^a   
\end{equation*}
has no solution for $a \ge 1$.    
\end{lemma}

\begin{proof}
The proof follows exactly the same argument as that of Lemma~\ref{lemma1}. Reducing the equation modulo $p$ gives $p=3$, and hence $q=7$. 
Next, reducing this equation modulo 7, we obtain $3^{2a+1}\equiv1\pmod7$. Since the order of $3$ modulo $7$ is $6$, it follows that $6\mid(2a+1)$, which is impossible.     
\end{proof}
Hence the equation~(\ref{equation1}) has no solution for case A. Note that if $a=0$ then $x=1$ and $y=0$. 
The equation $ p^{2a+1}-1=2q^a$ becomes $p-1=2$ and and hence $p=3$. 
Therefore, the  equation~(\ref{equation1})  has the solution $(p,x,y,z)=(3,1,0,2)$. Since $y$ is assumed to be positive in Theorem~\ref{A}, this case is excluded from the proof Lemma~\ref{lemma2}.

\subsection{Case B: x even, y odd:} In this case, we write $x=2a$ and $y=2b+1$, where $a \ge1 $ and $b \ge 0$.  Then
\begin{equation*}
q^{2b+1} = z^2-p^{2a} = (z-p^a)(z+p^a).    
\end{equation*}
We now consider two cases. If one of the factors is equal to $1$, we call it the \emph{trivial split}. 
Otherwise, both factors are greater than $1$, and we call it the \emph{non-trivial split}. We first discuss the trivial split.

In case of a trivial split, we have $z-p^a=1$ and $z+p^a=q^{2b+1}$. Subtracting the first equation from the second, we obtain 
\begin{equation} \label{equation3}
 q^{2b+1}-1=2p^a.   
\end{equation}
Since $x=2a$, $y=2b+1$, and $|x-y|=1$, there are two sub-cases:
\begin{enumerate}
\item[(i)] $y=x-1$, which is equivalent to $a=b+1$;
\item[(ii)] $y=x+1$, which is equivalent to $a=b$.
\end{enumerate}
For sub-case (i), the equation~(\ref{equation3}) becomes $(2p+1)^{2b+1}-1=2p^{\,b+1}$. The equation admits a solution for $b=0$ corresponding to $(x,y,z)=(2,1,p+1)$. The following proposition shows that this is the only solution.  
\begin{proposition}\label{P2.5}
\normalfont For every odd prime $p$, the equation $(2p+1)^{2b+1}-1=2p^{\,b+1}$
has a unique solution $b=0$,  giving $(x,y,z)=(2,1,p+1)$.
\end{proposition}

\begin{proof}
If $b=0$, then the equation becomes $(2p+1)-1=2p$, which holds. This gives the solution $(x,y,z)=(2,1,p+1)$ because $p^2+(2p+1)=(p+1)^2$. 
Now suppose that $b\ge1$. Since $q=2p+1\equiv1\pmod p$ and $v_p(q-1)=v_p(2p)=1$, from corollary~(\ref{corollary1}), we obtain $v_p(q^{2b+1}-1)=1+v_p(2b+1).$ 
By comparing the $p$-adic valuations of both sides of $q^{2b+1}-1=2p^{\,b+1}$, we obtain $1+v_p(2b+1)=b+1$, or equivalently, $v_p(2b+1)=b$. 
Hence $p^b\mid(2b+1)$ which implies $p^b\le2b+1$. If $b=1$, then $p\le3$. 
Since $p$ is an odd prime, we must have $p=3$, but $7^3-1=342\neq18=2\cdot3^2$, which is impossible. 
Next, for $b \ge 2$,  we claim that $p^b>2b+1$. This follows by induction on $b$. The statement is true for $b=2$, since $p^2\ge9>5$. 
Assume that $p^b>2b+1$ for some $b\ge2$. Then
\begin{equation*}
    p^{b+1}\ge3p^b>3(2b+1)>2(b+1)+1.
\end{equation*}
Therefore the equation $(2p+1)^{2b+1}-1=2p^{\,b+1}$ has no solution for $b \geq 1$ and an unique solution at $b=0$.    
\end{proof}

We now consider sub-case (ii), where $a=b$. In this case, equation~(\ref{equation3}) becomes $(2p+1)^{2b+1}-1=2p^{b}$. 
This sub-case also leads to no solution of equation~(\ref{equation3}).

\begin{lemma}
\normalfont For every odd prime $p$, the equation $q^{2b+1}-1 = 2p^b$ has no solution for any integer $b\ge0$.
\end{lemma}

\begin{proof}
If $b=0$, then the equation becomes $2p=2$, which is impossible. Now suppose that $b\ge1$. Since $2b+1\ge3$, we have  
\begin{equation*}
(2p+1)^{2b+1}-1\ge(2p+1)^3-1=8p^3+12p^2+6p.    
\end{equation*}
Moreover,
\begin{equation*}
    8p^3+12p^2+6p-2p^b = p\left(8p^2+12p+6-2p^{\,b-1}\right).
\end{equation*}
If $b=1$, then $8p^2+12p+6-2=8p^2+12p+4>0$. Now, let $b\ge2$. Since
\begin{equation*}
(2p+1)^{2b+1}\ge(2p+1)^3(2p+1)^{2b-2} \,\, \text{and} \,\,  (2p+1)^2>p,   
\end{equation*}
it follows that $(2p+1)^{2b+1}>(2p+1)^3p^{b-1} > 2p \cdot p^{b-1} = 2 p^b $. Therefore $(2p+1)^{2b+1}-1>2p^b$ and hence the equation has no solution.
\end{proof}
The trivial split produces only the solution $(x,y,z)=(2,1,p+1)$. Next, we consider non-trivial split case. 
In this case, both factors in $q^{2b+1}=(z-p^a)(z+p^a)$ are greater than $1$. Since $\gcd(z-p^a,z+p^a)=1$, every prime power dividing $q^{2b+1}$ must be assigned entirely to one of the two factors.
Thus there exist coprime integers $r$ and $s$ satisfying $q=rs$, with $1<r<s$,
such that $z-p^a=r^{\,2b+1}$ and $z+p^a=s^{\,2b+1}$. 
The difference of these two equations gives $s^{\,2b+1}-r^{\,2b+1}=2p^a$. 
Now, from the factorization of $s^{\,2b+1}-r^{\,2b+1}$, we obtain 
\begin{equation}\label{eq4}
(s-r)\sum_{i=0}^{2b}s^{\,2b-i}r^i=2p^a.
\end{equation}
For our convenience, let $\displaystyle S=\sum_{i=0}^{2b}s^{\,2b-i}r^i$. Then equation~\eqref{eq4} becomes $(s-r)S=2p^a$. Next, we want determine the possible values of $s - r$.
\begin{lemma}
\normalfont The gcd of $s-r$ and $S$ is $1$.    
\end{lemma}

\begin{proof}
Let $d=\gcd(s-r,S)$. Since $s\equiv r \pmod{s-r}$,  we have $S\equiv\sum_{i=0}^{2b}r^{2b} =(2b+1)r^{2b} \pmod{s-r}$. 
Hence $d\mid(2b+1)r^{2b}$. Since $\gcd(r,s-r)=\gcd(r,s)=1$, it follows that
$\gcd(r^{2b},s-r)=1$. Therefore, $d\mid(2b+1)$. 

As $d \mid (s-r) $ and $d\mid S$, we have $d^{\,2} \mid (s-r) S = 2 p^a$. In addition, $d$ is odd, and it follows that $d^{\,2}\mid p^a$.
Therefore, $d=p^t$ for some integer $t\ge0$. If $d>1$, then $d\ge p$, and since $d\mid(s-r)$, we obtain $s-r\ge d\ge p$.  But $r \ge 3$, as $r \mid q$ and $r >1$, so we have
\begin{equation*}
s=\frac{2p+1}{r}\le\frac{2p+1}{3}<p.
\end{equation*} 
Therefore $s-r<s<p\le d\le s-r$, which is impossible. Hence $d=1$.
\end{proof}

Since $\gcd(s-r,S)=1$ and $(s-r)S=2p^a$, therefore $2p^a$ has a co-prime factorization. 
If $b \ge 1$, then $S= \displaystyle \sum_{i=0}^{2b}s^{2b-i}r^i\ge 2b+1\ge3$. 
Hence the only possible coprime factorizations are $(s-r,S)=(1,2p^a)$ or $(s-r,S)=(2,p^a)$. Therefore $s-r\in\{1,2\}$. 
Now, suppose $s-r=1$, equivalently, $s=r+1$. since $rs=2p+1$, we obtain $p = \frac{r(r+1)-1}{2}$.
However, $r(r+1)$ is even, so $r(r+1)-1$ is odd. Hence $\frac{r(r+1)-1}{2}$ is not an integer, which is a contradiction. Therefore, $s-r\neq1$ and hence $s-r=2$.

Since both r and s are divisors of an odd integer and $s-r=2$,  we can write $s=2k-1$ and $r=2k+1$ for some $k \geq 2$. From $rs=2p+1$, we obtain $(2k-1)(2k+1)=2p+1$, which simplifies to $4k^2-1=2p+1$. 
Hence $p=2k^2-1$. Thus every non-trivial solution must satisfy $2p+1=(2k-1)(2k+1)$, where $p=2k^2-1$ is prime.

\begin{remark}\label{R2.8}
In the sub-case $a=b$ (equivalently, $y=x+1$), the sum $S= \displaystyle \sum_{i=0}^{2b}s^{2b-i}r^i$ contains the term $s^br^b=(rs)^b=(2p+1)^b>p^b$. Hence $S>(2p+1)^b>p^b=p^a$,
which is impossible. Therefore the sub-case $a=b$ has no solution.
\end{remark}

Therefore, the only remaining subcase is $a=b+1$. Now, if we substitute $r=2k-1$, $s=2k+1$ and $p=2k^2-1$ into the equation $s^{\,2b+1}-r^{\,2b+1}=2p^a$, then we obtain
\begin{equation*}
(2k+1)^{\,2b+1}-(2k-1)^{\,2b+1} = 2(2k^2-1)^{\,b+1}	
\end{equation*}
\begin{comment}
\begin{proposition}\label{P2.9}
	\normalfont
	Let $k\ge2$ and $b\ge1$ be integers. If
	\begin{equation*}
		(2k+1)^{2b+1}-(2k-1)^{2b+1} = 2(2k^2-1)^{\,b+1},
	\end{equation*}
	then 
	\begin{equation*}	
		2k^2-1 > (2b+1)2^b.
	\end{equation*}
\end{proposition}
\end{comment}
\section{Proof of the main result}
Now, using the results proved above, we now obtain the main theorem. The proof is divided into the trivial and non-trivial factorization cases, giving parts (i) and (ii), respectively.

\begin{theorem}\label{T3.1}
\normalfont Let $p$ be an odd prime and suppose that
\begin{equation*}
	p^x + \left ( 2p +1 \right)^y = z^2,
\end{equation*}
where $x$, $y$, $z$ are positive integers with $|x-y|=1$. Then exactly one of the following holds.
\begin{enumerate}
	\item[(i)] The solution is $(x, y, z) = (2, 1, p + 1)$.
	\item[(ii)] The integer $2p + 1$ is composite and admits the coprime factorization 
	\begin{equation*}
		2p + 1 = (2k - 1)(2k + 1),
	\end{equation*}
	where $p=2k^2-1$ and $(2k + 1)^{2b+1} - (2k - 1)^{2b+1} = 2(2k^2 - 1)^a$, with $(2b+1)2^b<2k^2-1$. 
	Consequently, $ b< \left\lceil \log_2(2k^2-1) \right\rceil$ leaving only a finite logarithmic range of exponents for each fixed prime. .
\end{enumerate}

\end{theorem}

\begin{proof}
	
From the identity $\displaystyle X^n-Y^n=(X-Y)\sum_{i=0}^{n-1}X^{n-1-i}Y^i$,	we have 
\begin{equation*}
	(2k+1)^{2b+1}-(2k-1)^{2b+1} = 2\sum_{i=0}^{2b}(2k+1)^{2b-i}(2k-1)^i.
\end{equation*}
Since $(2k+1)^{2b+1}-(2k-1)^{2b+1} = 2(2k^2-1)^{\,b+1}$, it follows that
\begin{equation}\label{e5}
	\sum_{i=0}^{2b}(2k+1)^{2b-i}(2k-1)^i = (2k^2-1)^{\,b+1}. 
\end{equation}
For each $0\le i\le b-1$, pair the $i$-th and $(2b-i)$-th terms.  By the arithmetic geometric mean inequality, we obtain
\begin{equation*}
	(2k+1)^{2b-i}(2k-1)^i + (2k+1)^i(2k-1)^{2b-i} \ge 2(2k+1)^b(2k-1)^b.
\end{equation*}
Since there are total $b$ such pairs along with the middle term $(2k+1)^b(2k-1)^b$, we obtain
\begin{equation*}
	\sum_{i=0}^{2b}(2k+1)^{2b-i}(2k-1)^i \ge (2b+1)(2k+1)^b(2k-1)^b.
\end{equation*}
Now combining this with equation~(\ref{e5}) gives $(2k^2-1)^{\,b+1}\ge(2b+1)(4k^2-1)^b$. Dividing both sides by $(2k^2-1)^b>0$, we obtain
\begin{equation*}
	2k^2-1 \ge	(2b+1)	\left(\frac{4k^2-1}{2k^2-1}\right)^b.
\end{equation*}
Since
\begin{equation*}
	\frac{4k^2-1}{2k^2-1} = 2+\frac{1}{2k^2-1}>2,
\end{equation*}
it follows that
\begin{equation}\label{e6}
	2k^2-1	>	(2b+1)2^b.
\end{equation}	
By Proposition~\ref{P2.5}, the equation has the solution $(x,y,z)=(2,1,p+1)$. If this is not the case, then Remark~\ref{R2.8} shows that $a=b$ is not possible. Hence $a=b+1$. 
Moreover, from equation~(\ref{e6}), we get $2k^2-1>(2b+1)2^b$. Now after taking logarithms, we obtain $b<\log_2(2k^2-1)$. 
Since $b$ is an integer, therefore
\begin{equation*}
	b<\left\lceil\log_2(2k^2-1)\right\rceil.
\end{equation*}
Thus, for each fixed $k$, only finitely many values of $b$ need to be considered.
\end{proof}
\begin{remark}
For $k=2$ and $b=1$, the auxiliary equation yields the solution	$p=7$, $(x,y,z)=(4,3,76)$, since $7^4+15^3=2401+3375=5776=76^2$. 
For odd integers $b\ge3$, no solutions were found in our computations. Although this suggests that the above solution may be unique, we do not have a general proof covering all such cases. 
Determining whether the auxiliary equation admits solutions beyond the example $p=7$ remains an interesting open problem.
\end{remark}

%\section*{Disclosure of Interest}

%No potential conflict of interest was reported by the author.


%\section*{Funding}

%No funding was received.

\bibliographystyle{amsplain}
\bibliography{mybib}


\end{document}


